\chapter{Pierwsze spostrzeżenia przed badaniem właściwym}

\section{Inwentaryzacja aplikacji po resecie}

\textbf{Cel:} uzyskanie stanu wyjściowego urządzeń po przywróceniu ustawień fabrycznych oraz zebranie metadanych o aplikacjach zainstalowanych domyślnie, ze szczególnym uwzględnieniem możliwości kontroli wykorzystania danych sieciowych w tle.

\textbf{Sprzęt:} dwa telefony Samsung Galaxy S7 (po resecie do ustawień fabrycznych).

Po oczyszczeniu zrzutów z urządzeń uzyskano spójny zbiór \textbf{53} aplikacji domyślnych na każdym z telefonów.

\textbf{Wyniki zbiorcze:}
\begin{table}[htbp]
  \centering
  \caption{Podsumowanie ustawień korzystania z danych sieciowych w tle (po resecie)}
  \label{tab:default-apps-summary}
  \begin{tabular}{lrr}
    \toprule
    \textbf{Kategoria} & \textbf{Liczba} & \textbf{Udział} \\
    \midrule
    Łączna liczba aplikacji      & 53 & 100\% \\
    Można zablokować dane w tle  & 44 & 83{,}0\% \\
    Nie można zablokować         & 9  & 17{,}0\% \\
    Domyślnie aktywny transfer   & 53 & 100\% \\
    Domyślnie zablokowany        & 0  & 0\% \\
    \bottomrule
  \end{tabular}
  \captionsetup{justification=raggedright,singlelinecheck=false}
  \caption*{\footnotesize
  \textbf{Legenda:} 
  „Można zablokować dane w tle” – aplikacje, dla których system Android udostępnia opcję ograniczenia transmisji danych w tle; 
  „Nie można zablokować” – aplikacje i komponenty systemowe zachowujące dostęp do sieci w tle niezależnie od użytkownika; 
  „Domyślnie aktywny transfer/zablokowany” – stan tuż po resecie fabrycznym; 
  „Udział” – procent względem wszystkich zidentyfikowanych aplikacji ($N=53$).
  }
\end{table}

\textbf{Aplikacje z nieograniczonym dostępem do danych w tle:}
\begin{longtable}{@{}l@{}}
\caption{Lista aplikacji, dla których nie można zablokować danych sieciowych w tle (po resecie)}\label{tab:default-apps-non-disablable}\\
\toprule
\textbf{Nazwa aplikacji} \\
\midrule
\endfirsthead
\toprule
\textbf{Nazwa aplikacji} \\
\midrule
\endhead
\bottomrule
\endfoot
Konserwacja urządzenia \\
Samsung Cloud \\
Samsung Members \\
Sim AT \\
SmartThings \\
Udostępnianie łącza \\
Ustawienia \\
Workspace \\
% (W CSV znajdowała się duplikacja „Sim AT”; po deduplikacji pozostaje jedna pozycja.)
\end{longtable}

\textbf{Wniosek:} stan „po resecie” obejmuje wyłącznie aplikacje \emph{domyślnie aktywne}, z których zdecydowaną większość (83\%) można ręcznie ograniczyć pod kątem wykorzystania danych w tle. Około \emph{17\%} stanowią jednak komponenty, które zachowują pełny dostęp do sieci w tle i nie podlegają tej kontroli.

\medskip
\noindent\emph{Uwaga replikacyjna:} powyższe wyniki pochodzą z pojedynczego zrzutu po resecie na tym samym systemie operacyjnym w obu urządzeniach.

\vspace{1ex}

\section{Przechwytywanie radiowe Wiresharkiem w trybie monitorowania (niepowodzenie)}
\textbf{Cel:} pasywna obserwacja ruchu bez konfigurowania proxy na urządzeniu.

\textbf{Sprzęt:} laptop z Linuksem; karta \emph{Wi-Fi} w trybie monitorowania; dwa Galaxy S7.

\textbf{Metoda:} nasłuch (\emph{monitor mode}) podczas korzystania z \emph{Wi-Fi} zabezpieczonego (\emph{WPA2}), próba późniejszej dekrypcji ramek i analizy ruchu aplikacji.

\textbf{Wynik:} \textit{niepowodzenie metody} w kontekście pełnej analizy treści: kluczowy ruch pozostaje szyfrowany (\emph{WPA2}), brak materiału do bezstratnej dekrypcji (m.in. konieczność złapania handshake i posiadania parametrów sieci). Metoda okazała się niewspółmiernie pracochłonna i ograniczona względem celów projektu; zrezygnowano na rzecz podejścia \emph{proxy}/\emph{pcap}.

\textbf{Wnioski:} przy analizie aplikacji mobilnych efektywniejsze jest środowisko kontrolowane (\emph{AP} + \emph{proxy}/\emph{pcap}) niż sniffing radiowy na produkcyjnej sieci.
% ----------
\section{\emph{MITM} przez \emph{proxy}}

\textbf{Cel:} przechwytywanie i inspekcja ruchu \emph{HTTP(S)} generowanego przez badane aplikacje.

\textbf{Narzędzia:} \textit{mitmproxy}, \textit{tcpdump}.

\textbf{Sprzęt i sieć:} laptop z systemem Debian skonfigurowany jako punkt dostępowy (\emph{hostapd} + \emph{dnsmasq}) oraz dwa telefony Samsung Galaxy S7 po resecie do ustawień fabrycznych.

\textbf{Metoda:}
\begin{enumerate}
  \item Konfiguracja \emph{AP} i \emph{routingu}; uruchomienie \textit{mitmproxy} w trybie przechwytującym oraz równoległe rejestrowanie pakietów za pomocą \textit{tcpdump}.
  \item Instalacja certyfikatu użytkownika \textit{mitmproxy} na urządzeniach oraz próby uruchomienia i obsługi wybranych aplikacji.
\end{enumerate}

\textbf{Wynik:} 
Przy podejściu pierwszym aplikacje odmawiały działania, wykrywając obecność \emph{proxy} i blokując komunikację. Próba drugiego wariantu (instalacja certyfikatu i obejście \emph{pinningu}) również zakończyła się niepowodzeniem: telefony odrzucały certyfikat podczas pracy, a dodatkowo podczas próby uzyskania uprawnień \emph{root} w celu głębszej ingerencji doszło do uszkodzenia jednego z urządzeń. Sytuacja ta wymusiła zmianę metodologii i poszukiwanie alternatywnych rozwiązań.

\textbf{Wnioski:} samo zastosowanie \emph{proxy MITM} daje bardzo ograniczoną obserwowalność wobec aplikacji stosujących \emph{pinning} i/lub protokół \emph{QUIC}. Niezbędne są inne podejścia, np.\ instrumentacja, przechwytywanie ruchu na poziomie systemowym (\emph{pcap}) lub konwersja ruchu \emph{HTTP/3} na \emph{HTTP/2}.

\section{\emph{MITM} z instrumentacją (\emph{Frida}) wobec \emph{pinningu} — próba nieudana}

\textbf{Cel:} obejście mechanizmu \emph{certificate pinning} i przywrócenie możliwości inspekcji treści żądań.

\textbf{Narzędzia:} \textit{Frida} (\textit{frida-server} oraz gotowe skrypty \emph{hookujące} interfejsy \emph{API} związane z walidacją certyfikatów).

\textbf{Metoda:} uruchomiono \textit{frida-server} na urządzeniu i wstrzyknięto skrypty mające na celu przechwycenie wywołań metod weryfikujących certyfikaty oraz ich dynamiczne obejście.

\textbf{Wynik:} próba zakończyła się \textit{niepowodzeniem}. Wstrzyknięte skrypty nie zadziałały w przypadku analizowanej aplikacji. Prawdopodobnymi przyczynami były: brak uprawnień \emph{root}, aktywne mechanizmy \emph{SELinux} oraz natywne implementacje weryfikacji certyfikatów (np.\ niestandardowe \texttt{TrustManager}), których nie obejmowały standardowe \emph{hooki}.

\textbf{Wnioski:} w wielu nowoczesnych aplikacjach obejście \emph{pinningu} bez dostępu do uprawnień \emph{root} jest nierealne lub bardzo trudne. Alternatywami pozostają uruchamianie aplikacji w środowisku z pełnym dostępem administracyjnym (emulator, urządzenie z \emph{root}) lub analiza ograniczona do metadanych transportowych (hosty, porty, czasy, wolumen danych), bez wglądu w same treści komunikatów.


% % ----------
% \section{Rejestracja ruchu po stronie urządzenia (PCAP/PCAPNG) — przykład Instagram (sukces)}
% \textbf{Cel:} charakterystyka ruchu bez naruszania zabezpieczeń aplikacji (rejestracja na urządzeniu).

% \textbf{Metoda:} rejestracja ruchu i analiza strumieni; w tej próbie dominował HTTP/3/QUIC (UDP/443) z kanałem utrzymania sesji po TCP/5222, a ruch HTTP/1.1/HTTPS był śladowy.
% % Liczby i wykresy trzymamy w rozdziale z wynikami – tu tylko krótkie streszczenie

% \textbf{Wynik:} metoda skuteczna do obserwacji protokołów, kierunków i aktywności w tle; semantyka HTTP ograniczona przez QUIC/pinning. 
% \textbf{Artefakty:} patrz rozdział z eksperymentem głównym (plik .pcapng i agregaty).
